\documentclass[10pt]{article}

\usepackage[T1]{fontenc}
\usepackage{lmodern}
% The bundled offline compiler lacks only ec-lmr6.tfm; use the adjacent
% optical size for 6pt subscripts while retaining the standard arXiv font stack.
\DeclareFontShape{T1}{lmr}{m}{n}{%
  <-5.5> ec-lmr5
  <5.5-7.5> ec-lmr7
  <7.5-8.5> ec-lmr8
  <8.5-9.5> ec-lmr9
  <9.5-11> ec-lmr10
  <11-15> ec-lmr12
  <15-> ec-lmr17}{}
\usepackage{microtype}
\usepackage[margin=0.78in]{geometry}
\usepackage{amsmath,amssymb}
\usepackage{booktabs}
\usepackage{tabularx}
\usepackage{array}
\usepackage{enumitem}
\usepackage{xcolor}
\usepackage{graphicx}
\usepackage{caption}
\usepackage{etoolbox}
\usepackage[numbers,sort&compress]{natbib}
\usepackage[hidelinks]{hyperref}
\usepackage[nameinlink,noabbrev]{cleveref}
\crefname{table}{Table}{Tables}
\crefname{figure}{Figure}{Figures}
\crefname{section}{Section}{Sections}

% Single source of truth for all empirical quantities in the manuscript.
% Update this file only after a validated experiment and statistics export.
\newcommand{\ResultPending}[1]{\textsc{Pending}}
\newcommand{\NotApplicable}{\textsc{N/A}}
\newcommand{\MeasuredResult}[1]{#1}
\newcommand{\ProjectedTarget}[1]{#1}

% Evidence switch. Keep false until a matched SafeDesk run is validated.
\newif\ifSafeResultsAvailable
\SafeResultsAvailablefalse

% Dataset and protocol constants.
\newcommand{\AWTaskCount}{417}
\newcommand{\AWScenarioCount}{139}
\newcommand{\AWAppCount}{9}
\newcommand{\AWAPICount}{457}
\newcommand{\AWAPILimit}{20}
\newcommand{\AWMaxTurns}{50}
\newcommand{\AWShardCount}{4}
\newcommand{\AWBootstrapSamples}{10{,}000}
\newcommand{\ModelTemperature}{0}
\newcommand{\TauDomainCount}{3}
\newcommand{\TauDiagnosticTaskCount}{278}
\newcommand{\TauTrialsPerTask}{4}
\newcommand{\TauPackageVersion}{1.0.0}
\newcommand{\TauBenchmarkRevision}{\texttt{1901a301961cbbe3fd11f3e84a2a376530c759e3}}
\newcommand{\TauTaskListHash}{\texttt{755c997d56c0e50275aedc823b27ffe54b2a967f12b20dea30df25a02f1b217e}}

% Local implementation-validation record; this is not a benchmark result.
\newcommand{\LocalUnitTestPassed}{142}
\newcommand{\LocalUnitTestSkipped}{1}
\newcommand{\AppWorldCatalogToolCount}{457}
\newcommand{\PrototypeValidationTaskCount}{20}
\newcommand{\PrototypeValidationRepetitionCount}{20}
\newcommand{\PrototypeValidationPassed}{20}
\newcommand{\PrototypeValidationCrashes}{0}

% AppWorld measured DeepSeek V4 Flash baseline.
\newcommand{\AWBaseSuccessCount}{174}
\newcommand{\AWBaseTGCValue}{41.73}
\newcommand{\AWBaseTGC}{\AWBaseTGCValue\%}
\newcommand{\AWBaseTGCClusterCILow}{34.53}
\newcommand{\AWBaseTGCClusterCIHigh}{49.16}
\newcommand{\AWBaseScenarioSuccessCount}{41}
\newcommand{\AWBaseSGCValue}{29.50}
\newcommand{\AWBaseSGC}{\AWBaseSGCValue\%}
\newcommand{\AWBaseEvaluatorPassed}{2{,}218}
\newcommand{\AWBaseEvaluatorTotal}{3{,}348}
\newcommand{\AWBaseEvaluatorPassValue}{66.25}
\newcommand{\AWBaseEvaluatorPass}{\AWBaseEvaluatorPassValue\%}

\newcommand{\AWBaseCompletionTaskCount}{327}
\newcommand{\AWBaseProposedFalseCompletionCount}{153}
\newcommand{\AWBaseProposedFalseCompletionValue}{46.79}
\newcommand{\AWBaseProposedFalseCompletion}{\AWBaseProposedFalseCompletionValue\%}
\newcommand{\AWBaseSystemAcceptedCompletionCount}{327}
\newcommand{\AWBaseSystemAcceptedFailureCount}{153}
\newcommand{\AWBaseSystemAcceptedFailureValue}{46.79}
\newcommand{\AWBaseSystemAcceptedFailure}{\AWBaseSystemAcceptedFailureValue\%}
\newcommand{\AWBaseNoCompletionCount}{90}
\newcommand{\AWBaseNoCompletionValue}{21.58}
\newcommand{\AWBaseNoCompletion}{\AWBaseNoCompletionValue\%}
\newcommand{\AWBaseMaxTurnCount}{91}
\newcommand{\AWBaseMaxTurnValue}{21.82}
\newcommand{\AWBaseMaxTurn}{\AWBaseMaxTurnValue\%}
\newcommand{\AWBaseFailureCount}{243}
\newcommand{\AWBaseFailureWithPassedTestCount}{237}

\newcommand{\AWBaseProposedCallCount}{15{,}066}
\newcommand{\AWBaseExecutedCallCount}{14{,}631}
\newcommand{\AWBaseReadCallCount}{12{,}795}
\newcommand{\AWBaseWriteCallCount}{1{,}836}
\newcommand{\AWBaseInvalidCount}{559}
\newcommand{\AWBaseInvalidValue}{3.71}
\newcommand{\AWBaseInvalid}{\AWBaseInvalidValue\%}
\newcommand{\AWBaseOutOfSchemaCount}{435}
\newcommand{\AWBaseOutOfSchemaValue}{2.89}
\newcommand{\AWBaseOutOfSchema}{\AWBaseOutOfSchemaValue\%}
\newcommand{\AWBaseDuplicateCallCount}{987}
\newcommand{\AWBaseDuplicateCallValue}{6.75}
\newcommand{\AWBaseDuplicateCall}{\AWBaseDuplicateCallValue\%}
\newcommand{\AWBaseDuplicateWriteCount}{128}
\newcommand{\AWBaseDuplicateWriteValue}{6.97}
\newcommand{\AWBaseDuplicateWrite}{\AWBaseDuplicateWriteValue\%}
\newcommand{\AWBaseOutOfSchemaAffectedTaskCount}{217}
\newcommand{\AWBaseOutOfSchemaAffectedTaskValue}{52.04}
\newcommand{\AWBaseOutOfSchemaAffectedTaskRate}{\AWBaseOutOfSchemaAffectedTaskValue\%}
\newcommand{\AWBaseDuplicateAffectedTaskCount}{179}
\newcommand{\AWBaseDuplicateAffectedTaskValue}{42.93}
\newcommand{\AWBaseDuplicateAffectedTaskRate}{\AWBaseDuplicateAffectedTaskValue\%}
\newcommand{\AWBaseDuplicateWriteAffectedTaskCount}{40}
\newcommand{\AWBaseDuplicateWriteAffectedTaskValue}{9.59}
\newcommand{\AWBaseDuplicateWriteAffectedTaskRate}{\AWBaseDuplicateWriteAffectedTaskValue\%}

\newcommand{\AWBaseTotalTokens}{201{,}973{,}864}
\newcommand{\AWBasePredictorTokens}{3{,}198{,}312}
\newcommand{\AWBaseAgentTokens}{198{,}775{,}552}
\newcommand{\AWBaseMeanTokens}{484{,}350}
\newcommand{\AWBaseMedianTokens}{136{,}548}
\newcommand{\AWBasePNinetyFiveTokens}{2{,}056{,}783}
\newcommand{\AWBaseMaxTokens}{7{,}128{,}790}
\newcommand{\AWBaseSuccessMeanTokens}{130{,}653}
\newcommand{\AWBaseSuccessMedianTokens}{86{,}875}
\newcommand{\AWBaseSuccessPNinetyFiveTokens}{307{,}502}
\newcommand{\AWBaseSuccessMaxTokens}{2{,}132{,}035}
\newcommand{\AWBaseFailureMeanTokens}{737{,}614}
\newcommand{\AWBaseFailureMedianTokens}{437{,}490}
\newcommand{\AWBaseFailurePNinetyFiveTokens}{2{,}848{,}263}
\newcommand{\AWBaseFailureMaxTokens}{7{,}128{,}790}
\newcommand{\AWBaseMeanTurns}{20.62}
\newcommand{\AWBaseFailureMeanTurns}{28.46}
\newcommand{\AWBaseMeanCalls}{35.09}
\newcommand{\AWBaseValidDurationCount}{409}
\newcommand{\AWBaseMeanDuration}{105.52}
\newcommand{\AWBaseMedianDuration}{37.29}
\newcommand{\AWBasePNinetyFiveDuration}{493.19}
\newcommand{\AWBaseMaxDuration}{1{,}776.09}
\newcommand{\AWBaseCostLow}{208.38}
\newcommand{\AWBaseCostHigh}{625.15}

% Baseline validity and provider identifiers.
\newcommand{\AWModelID}{\texttt{deepseek-v4-flash}}
\newcommand{\AWAPIEndpoint}{\texttt{https://api.deepseek.com/v1}}
\newcommand{\AWModelAccessDate}{2026-07-15--16}
\newcommand{\AWProviderVersion}{not recorded}
\newcommand{\AWGBKRepairCount}{11}
\newcommand{\AWGBKRepairSuccessCount}{4}
\newcommand{\AWFinalInfraErrorCount}{0}
\newcommand{\AWInvalidDurationCount}{8}
\newcommand{\AWTaskListHash}{97314d90d8600cba7c2bba8d28f9fee4a05ce934a9a7f0d2039278d316c77235}
\newcommand{\AWRepairListHash}{881818a44c46e10e5905041cbbdab9e0b42a8750a4cf14db364243b497121207}
\newcommand{\AWTaskListHashDisplay}{97314d90d8600cba\allowbreak{}7c2bba8d28f9fee4\allowbreak{}a05ce934a9a7f0d2\allowbreak{}039278d316c77235}
\newcommand{\AWRepairListHashDisplay}{881818a44c46e10e\allowbreak{}5905041cbbdab9e0\allowbreak{}b42a8750a4cf14db\allowbreak{}364243b497121207}

% Baseline difficulty strata and scenario-clustered bootstrap intervals.
\newcommand{\AWDiffOneCount}{72}
\newcommand{\AWDiffOneSuccess}{51}
\newcommand{\AWBaseDiffOneTGCValue}{70.83}
\newcommand{\AWBaseDiffOneCILow}{54.17}
\newcommand{\AWBaseDiffOneCIHigh}{84.72}
\newcommand{\AWDiffTwoCount}{150}
\newcommand{\AWDiffTwoSuccess}{60}
\newcommand{\AWBaseDiffTwoTGCValue}{40.00}
\newcommand{\AWBaseDiffTwoCILow}{28.00}
\newcommand{\AWBaseDiffTwoCIHigh}{52.67}
\newcommand{\AWDiffThreeCount}{195}
\newcommand{\AWDiffThreeSuccess}{63}
\newcommand{\AWBaseDiffThreeTGCValue}{32.31}
\newcommand{\AWBaseDiffThreeCILow}{23.08}
\newcommand{\AWBaseDiffThreeCIHigh}{42.56}

% Matched AppWorld SafeDesk results. These remain pending until observed.
\newcommand{\AWSafeSuccessCount}{\ResultPending{}}
\newcommand{\AWSafeTGC}{\ResultPending{}}
\newcommand{\AWTGCIncrease}{\ResultPending{}}
\newcommand{\AWSafeScenarioSuccessCount}{\ResultPending{}}
\newcommand{\AWSafeSGC}{\ResultPending{}}
\newcommand{\AWSGCIncrease}{\ResultPending{}}
\newcommand{\AWSafeEvaluatorPass}{\ResultPending{}}
\newcommand{\AWEvaluatorPassChange}{\ResultPending{}}
\newcommand{\AWSafeNoCompletion}{\ResultPending{}}
\newcommand{\AWNoCompletionChange}{\ResultPending{}}
\newcommand{\AWSafeGateAcceptedFalseCompletion}{\ResultPending{}}
\newcommand{\AWGateAcceptedFalseCompletionChange}{\ResultPending{}}
\newcommand{\AWSafeMaxTurn}{\ResultPending{}}
\newcommand{\AWMaxTurnChange}{\ResultPending{}}
\newcommand{\AWSafeInvalid}{\ResultPending{}}
\newcommand{\AWInvalidChange}{\ResultPending{}}
\newcommand{\AWSafeOutOfSchema}{\ResultPending{}}
\newcommand{\AWOutOfSchemaChange}{\ResultPending{}}
\newcommand{\AWSafeDuplicateCall}{\ResultPending{}}
\newcommand{\AWDuplicateCallChange}{\ResultPending{}}
\newcommand{\AWSafeDuplicateWrite}{\ResultPending{}}
\newcommand{\AWDuplicateWriteChange}{\ResultPending{}}
\newcommand{\AWSafeMeanTokens}{\ResultPending{}}
\newcommand{\AWMeanTokenChange}{\ResultPending{}}
\newcommand{\AWSafeMedianTokens}{\ResultPending{}}
\newcommand{\AWSafePNinetyFiveTokens}{\ResultPending{}}
\newcommand{\AWSafeMaxTokens}{\ResultPending{}}
\newcommand{\AWSafeCost}{\ResultPending{}}
\newcommand{\AWSafeDuration}{\ResultPending{}}

% Completion Gate task-level and attempt-level fields. Never merge these scopes.
\newcommand{\AWSafeCompletionAttemptCount}{\ResultPending{}}
\newcommand{\AWSafeCompletionTaskCount}{\ResultPending{}}
\newcommand{\AWSafeBlockedAttemptCount}{\ResultPending{}}
\newcommand{\AWSafeBlockedTaskCount}{\ResultPending{}}
\newcommand{\AWSafeBlockedThenSuccessCount}{\ResultPending{}}
\newcommand{\AWSafeBlockedThenFailureCount}{\ResultPending{}}
\newcommand{\AWSafeGateAcceptedTaskCount}{\ResultPending{}}
\newcommand{\AWSafeGateAcceptedFailureCount}{\ResultPending{}}
\newcommand{\AWSafeSystemAcceptedFailureRate}{\ResultPending{}}
\newcommand{\AWSafeModelProposedFailureRate}{\ResultPending{}}
\newcommand{\AWSafeFalseBlockCount}{\ResultPending{}}
\newcommand{\AWSafeFalseBlockRate}{\ResultPending{}}
\newcommand{\AWSafeFalseBlockAttemptCount}{\ResultPending{}}
\newcommand{\AWSafeFalseBlockAttemptRate}{\ResultPending{}}
\newcommand{\AWSafeFalseBlockAffectedTaskCount}{\ResultPending{}}
\newcommand{\AWSafeAverageBlocksPerRescuedTask}{\ResultPending{}}
\newcommand{\AWSafeGateRescueRate}{\ResultPending{}}
\newcommand{\AWSafeRecoveryAttemptCount}{\ResultPending{}}
\newcommand{\AWSafeRecoveryAttemptSuccessRate}{\ResultPending{}}
\newcommand{\AWSafeTaskSuccessAfterRecoveryRate}{\ResultPending{}}

% Formal paired statistics; populated only by the frozen statistics script.
\newcommand{\AWPairedBothSuccess}{\ResultPending{}}
\newcommand{\AWPairedBaseOnlySuccess}{\ResultPending{}}
\newcommand{\AWPairedSafeOnlySuccess}{\ResultPending{}}
\newcommand{\AWPairedBothFailure}{\ResultPending{}}
\newcommand{\AWMcNemarPValue}{\ResultPending{}}
\newcommand{\AWPairedTGCCI}{\ResultPending{}}
\newcommand{\AWPairedTokenCI}{\ResultPending{}}

% Tau2 diagnostic single-run observations. These are not formal Pass^4 results.
\newcommand{\TauDiagAirlineTasks}{50}
\newcommand{\TauDiagAirlineSuccess}{32}
\newcommand{\TauDiagAirlineRate}{64.00\%}
\newcommand{\TauDiagRetailTasks}{114}
\newcommand{\TauDiagRetailSuccess}{91}
\newcommand{\TauDiagRetailRate}{79.82\%}
\newcommand{\TauDiagTelecomTasks}{114}
\newcommand{\TauDiagTelecomSuccess}{113}
\newcommand{\TauDiagTelecomRate}{99.12\%}
\newcommand{\TauAgentModelID}{\texttt{glm-5}}
\newcommand{\TauAgentEndpoint}{\texttt{https://open.bigmodel.cn/api/paas/v4/}}
\newcommand{\TauModelAccessDate}{2026-07-11--16}
\newcommand{\TauProviderVersion}{not recorded}

% Formal Tau2 matched results, requiring four fixed trials per task.
\newcommand{\TauBaseAirlinePassOne}{\ResultPending{}}
\newcommand{\TauSafeAirlinePassOne}{\ResultPending{}}
\newcommand{\TauBaseAirlinePassFour}{\ResultPending{}}
\newcommand{\TauSafeAirlinePassFour}{\ResultPending{}}
\newcommand{\TauBaseRetailPassOne}{\ResultPending{}}
\newcommand{\TauSafeRetailPassOne}{\ResultPending{}}
\newcommand{\TauBaseRetailPassFour}{\ResultPending{}}
\newcommand{\TauSafeRetailPassFour}{\ResultPending{}}
\newcommand{\TauBaseTelecomPassOne}{\ResultPending{}}
\newcommand{\TauSafeTelecomPassOne}{\ResultPending{}}
\newcommand{\TauBaseTelecomPassFour}{\ResultPending{}}
\newcommand{\TauSafeTelecomPassFour}{\ResultPending{}}
\newcommand{\TauBaseMacroPassOne}{\ResultPending{}}
\newcommand{\TauSafeMacroPassOne}{\ResultPending{}}
\newcommand{\TauPassOneIncrease}{\ResultPending{}}
\newcommand{\TauBaseMacroPassFour}{\ResultPending{}}
\newcommand{\TauSafeMacroPassFour}{\ResultPending{}}
\newcommand{\TauPassFourIncrease}{\ResultPending{}}

% Ablation results.
\newcommand{\AblFullTGC}{\ResultPending{}}
\newcommand{\AblWithoutSVTGC}{\ResultPending{}}
\newcommand{\AblWithoutTEGTGC}{\ResultPending{}}
\newcommand{\AblWithoutRCTGC}{\ResultPending{}}
\newcommand{\AblWithoutCMTGC}{\ResultPending{}}
\newcommand{\AblIncrementalBaseTGC}{\AWBaseTGC}
\newcommand{\AblIncrementalSVTGC}{\ResultPending{}}
\newcommand{\AblIncrementalTEGTGC}{\ResultPending{}}
\newcommand{\AblIncrementalRCTGC}{\ResultPending{}}
\newcommand{\AblIncrementalCMTGC}{\ResultPending{}}

% Development targets from a synthetic projection. These macros are forbidden
% in main.tex, tables, captions, abstract, and conclusion until observed.
\newcommand{\ProjectedAWSafeSuccessCount}{256}
\newcommand{\ProjectedAWSafeTGC}{61.39\%}
\newcommand{\ProjectedAWTGCIncrease}{+19.66 pp}
\newcommand{\ProjectedAWSafeScenarioSuccessCount}{65}
\newcommand{\ProjectedAWSafeSGC}{46.76\%}
\newcommand{\ProjectedAWSGCIncrease}{+17.27 pp}
\newcommand{\ProjectedAWSafeGateAcceptedFalseCompletion}{17.42\%}
\newcommand{\ProjectedAWSafeMeanTokens}{280{,}198}
\newcommand{\ProjectedTauPassOneIncrease}{+6.05 pp}
\newcommand{\ProjectedTauPassFourIncrease}{+12.14 pp}

% Fill these fields with real submission metadata before public release.
% They are intentionally empty rather than populated with invented identities.
\newcommand{\PaperAuthors}{}
\newcommand{\PaperAffiliation}{}
\newcommand{\PaperEmail}{}
\newcommand{\CodeRepository}{}
\newcommand{\AuthorContributions}{}
\newcommand{\ArtifactLicense}{}


\definecolor{safeblue}{HTML}{276FBF}
\definecolor{safegreen}{HTML}{2A9D8F}
\definecolor{safecoral}{HTML}{C8553D}
\definecolor{safegold}{HTML}{B77900}
\definecolor{safelight}{HTML}{F3F5F7}
\definecolor{safedark}{HTML}{263442}

\hypersetup{
  pdftitle={SafeDesk: Evidence-Grounded Runtime Governance for Long-Horizon Tool-Using Agents},
  pdfauthor={\PaperAuthors},
  pdfsubject={Runtime governance for reliable tool-using language-model agents},
  pdfkeywords={LLM agents, tool calling, runtime governance, verification, recovery, context management},
  colorlinks=true,
  linkcolor=safeblue,
  citecolor=safeblue,
  urlcolor=safeblue
}

\captionsetup{hypcap=false,font=small,labelfont=bf}
\newcolumntype{Y}{>{\raggedright\arraybackslash}X}
\newcolumntype{R}{>{\raggedleft\arraybackslash}X}
\newcommand{\safedesk}{SafeDesk}
\newcommand{\tgc}{\mathrm{TGC}}
\newcommand{\sgc}{\mathrm{SGC}}
\newcommand{\code}[1]{\texttt{#1}}
\setlength{\emergencystretch}{1.5em}
\setlist{leftmargin=*,itemsep=2pt,topsep=3pt}

\title{\safedesk: Evidence-Grounded Runtime Governance\\
for Long-Horizon Tool-Using Agents\\[0.45em]
\large Runtime Design, Prototype Tests, and a Proposed Evaluation Protocol}
\author{\PaperAuthors%
  \ifdefempty{\PaperAffiliation}{}{\\\small\PaperAffiliation}%
  \ifdefempty{\PaperEmail}{}{\\\small\texttt{\PaperEmail}}}
\date{\today}

\begin{document}
\maketitle
\ifdefempty{\CodeRepository}{}{\begin{center}\small Code: \url{\CodeRepository}\end{center}}

\begin{abstract}
Long-horizon tool-using agents fail through more than isolated reasoning errors:
goals drift, subgoals are omitted, invalid calls are proposed, side effects are
repeated, failures trigger blind retries, completion is asserted without evidence,
and growing histories bury critical constraints. We introduce \safedesk, a
model-interface-agnostic runtime governance layer that leaves model weights unchanged and
coordinates four modules: State \& Verification, Tool Execution Guard, Recovery
Controller, and Context Manager. The runtime separates model-proposed actions,
runtime-authorized execution and completion, and evaluator-confirmed outcomes.
It induces a task contract, tracks admissible evidence and write effects,
validates and schedules calls before execution, applies typed recovery after
classified failures, and assembles bounded model contexts without discarding hard
constraints. Completion is permitted only when required goals, effects, evidence,
and policy conditions are satisfied.
We implement a locally unit-tested runtime prototype, validate one deterministic
end-to-end mock scenario over \PrototypeValidationRepetitionCount{} repetitions,
and unit-test schema conversion against the AppWorld public API catalog. We also
specify a proposed, resource-capped protocol for AppWorld and $\tau^2$-bench. A
historical AppWorld diagnostic is retained only in
the appendix and is not a SafeDesk comparison. We make no claim here that
SafeDesk improves benchmark performance: that claim requires a newly constructed,
independently administered held-out evaluation with interleaved conditions and
recorded provider response metadata. The resulting artifact is a falsifiable
runtime design, a locally tested prototype, and a prespecified protocol for
studying whether governance complements stronger models and agentic training
without updating model weights.
\end{abstract}

\section{Introduction}

The central reliability gap in tool agents is not simply whether a language model
can generate a plausible API call. It is whether a multi-step execution lifecycle
preserves the user's goals, obeys policy, controls side effects, verifies state,
recovers from failures, and terminates on evidence. AppWorld exposes this gap with
\AWAppCount{} applications, \AWAPICount{} APIs, cross-application tasks, and
state-based evaluators \citep{trivedi2024appworld}. $\tau^2$-bench adds dual-control,
text-based half-duplex collaboration among a user, an agent, and domain tools
under business policies \citep{barres2025tau2,yao2024taubench}. Together they
distinguish syntactically valid tool use
from reliable task execution.

Consider a request to create a calendar event and notify its attendees. An agent
may omit the notification subgoal, submit a creation call with a missing field,
retry by creating a second event, trust a nominal success response without reading
the calendar back, and then declare completion. Each local action may look
plausible, yet the final environment contains a duplicate event and no notification.
A stronger prompt or a larger turn cap does not systematically track the omitted
goal, prevent the duplicate write, establish the postcondition, or ensure that the
final claim matches the persisted state.

We study the following question: \emph{Can a model-interface-agnostic runtime improve
long-horizon task completion, execution reliability, and resource efficiency
without updating model weights?} Training and agentic reinforcement learning can
improve the model's proposal distribution, but deployed systems still need
deterministic controls over what is executed, what is believed, and when completion
is allowed.

\safedesk{} addresses this execution layer through four composable modules. State
\& Verification maintains an explicit contract, evidence, and effects. Tool
Execution Guard validates and schedules calls before they reach the environment.
Recovery Controller classifies failures and applies bounded, typed recovery.
Context Manager reconstructs a compact model-facing view while preserving hard
invariants. The same runtime hook records every proposal, disposition, state
transition, and evaluator-facing artifact.

Our contributions are fourfold:
\begin{enumerate}
  \item \textbf{Runtime abstraction.} We separate task state, admissible evidence,
        side effects, failure state, and model context rather than treating message
        history as the only execution state.
  \item \textbf{Four governance modules.} We design composable controls for state
        verification, tool execution, typed recovery, and bounded context.
  \item \textbf{Runtime interface and prototype adapters.} We provide a
        framework-neutral interface design, Action IR, structured traces, a
        unit-tested AppWorld adapter prototype, and prototype benchmark adapters; their
        tested scope and untested end-to-end boundary are reported explicitly.
  \item \textbf{Falsifiable evaluation protocol.} We specify a future held-out,
        interleaved, fixed-resource comparison with leave-one-out ablations and
        analyses of completion, invalid calls, side effects, stagnation, and token cost.
\end{enumerate}

\section{Related Work}

\subsection{Reasoning and recovery agents}

ReAct interleaves reasoning and environmental action
\citep{yao2023react}; Reflexion uses verbal feedback across attempts
\citep{shinn2023reflexion}; and CRITIC revises outputs using tool-mediated critique
\citep{gou2023critic}. Voyager combines an automatic curriculum, a reusable skill
library, and environment-feedback-driven program refinement
\citep{wang2023voyager}. ReWOO separates planning from observations to reduce
repeated prompting and tool interleaving cost \citep{xu2023rewoo}. These methods
improve reasoning traces, feedback, plans, or reflection over tool results.
\safedesk{} is complementary: an external runtime can deterministically reject,
defer, verify, or recover a proposed action regardless of how the proposal was generated.

\subsection{Agent safety and guardrails}

ToolEmu uses an LM-emulated sandbox and evaluator to expose risky tool-agent
behavior \citep{ruan2024toolemu}. GuardAgent translates guard requests into
executable checks around another agent \citep{xiang2025guardagent}. Tool input and
output guardrails commonly validate schemas, filter sensitive fields, and enforce
local policy; durable workflow systems additionally rely on idempotent effects and
explicit retry boundaries. \safedesk{} extends the protected surface beyond a
single action: it governs task induction, dependency ordering, effect identity,
post-action verification, recovery, completion, and grounded response generation.

\subsection{Harness-level systems and evaluation}

Harness-Bench measures how realistic agent workflows change across harness
configurations while holding task environments, budgets, and evaluation protocols
comparable \citep{yao2026harnessbench}. AI Harness Engineering frames the harness
as a runtime substrate responsible for task specification, context, tool access,
state, observability, verification, permissions, and failure attribution
\citep{zhong2026aiharness}. These works directly motivate reporting the
model--harness configuration rather than attributing outcomes to the base model
alone. SafeDesk is neither a new benchmark nor a general taxonomy of harness
responsibilities. Its narrower contribution is a typed API-agent runtime design
for three enforceable properties: effect idempotency, evidence-grounded completion,
and authoritative-state isolation. The proposed AppWorld and $\tau^2$ protocols
are designed to test this runtime under matched budgets; they do not replace the
benchmark contribution of Harness-Bench.

\subsection{State and context management}

LangGraph persists graph state as ordered checkpoints and supports replay and
state updates through reducers \citep{langgraph2026persistence}. MemGPT manages
hierarchical memory to extend effective context \citep{packer2023memgpt}, while
InfiAgent externalizes persistent workspace state and reconstructs bounded context
for very long tasks \citep{yu2026infiagent}. Most closely related, StructAgent
organizes long-horizon agents around explicit task state, verifier-backed state
transitions, evidence-driven completion, and targeted recovery
\citep{wu2026structagent}. We do not claim priority for state-driven agents.
\safedesk{} instead specializes the abstraction for API/function-calling settings:
it adds an explicit Effect Ledger, schema repair and tool-boundary control,
dependency-aware scheduling, duplicate-write and idempotency checks, policy and
user-coordination gates for $\tau^2$-bench, and joint measurement of false
completion, invalid calls, side effects, and tokens.

\subsection{Training-based tool agents}

ToolLLM constructs large-scale API-use data and trains models with an API retriever
and search procedure \citep{qin2023toolllm}. Qwen3 supports thinking and
non-thinking modes and reports broad agent capability \citep{yang2025qwen3}.
Sequence-level agentic RL methods such as SeeUPO optimize multi-turn behavior
\citep{hu2026seeupo}, and the AppWorld leaderboard reports strong AgentRL systems
\citep{appworldleaderboard2026}. Runtime governance and agentic training are
complementary rather than mutually exclusive: training changes which actions are
proposed, while the runtime governs the execution lifecycle of those proposals.

The distinctions above are scoped claims about mechanisms, not a scorecard of
overall system quality. Direct empirical comparisons require matched tasks,
models, resources, and evaluators.

\section{Problem Formulation}

\subsection{State and task contract}

We use the notation summarized in \cref{tab:notation}.
\begin{center}
\begin{minipage}{0.82\linewidth}
\centering
\captionof{table}{Runtime notation.}
\label{tab:notation}
\small
\begin{tabular}{@{}cl@{}}
\toprule
Symbol & Meaning \\
\midrule
$W_t$ & environment world state \\
$\widehat{W}_t$ & runtime-observable projection of world state \\
$H_t$ & \safedesk{} task state \\
$E_t$ & Evidence Board \\
$L_t$ & Effect Ledger \\
$F_t$ & Failure State \\
$C_t$ & bounded Context Pack \\
$A_t$ & model-proposed actions \\
$O_t$ & tool observations \\
$\Delta W_t$ & observed environment-state change \\
\bottomrule
\end{tabular}
\end{minipage}
\end{center}

A task contract is
\begin{equation}
  \mathcal{T}=(G,D,\Gamma,P,V),
\end{equation}
where $G$ contains goals and subgoals, $D$ is the directed dependency relation on
those goals, $\Gamma$ contains user constraints and invariants, $P$ defines tool
permissions and business policies, and $V$ maps each required subgoal or intended
effect to admissible completion evidence. Unresolved fields $U_t$ are execution
state, not contractual requirements: they belong to $H_t$ and block relevant
transitions until resolved. At step $t$, the model receives $C_t$ and proposes
$A_t$. The runtime assigns each action a disposition, authorized tools return
$O_t$, and the State Reducer updates $(H_t,E_t,L_t,F_t)$ together with observed
$\Delta W_t$.

\subsection{Task Contract Induction}
\label{sec:contract-induction}

Contract induction maps a user instruction $u$, active tool schemas $\Sigma_t$,
and explicit policy $P$ to a typed proposal
$Q=(G,D,\Gamma,P,V,U)$: contractual goals, dependencies, constraints, policy,
evidence requirements, and unresolved execution fields. After deterministic
validation, the authoritative contract is $\mathcal T=(G,D,\Gamma,P,V)$ and $U$
is placed in $H_t$. The proposal conforms to a versioned JSON Schema. Each goal has
a stable identifier, required/optional status, arguments, dependencies, and a
verifier recipe; each unresolved field records its source clause and blocking scope.

\begin{center}
\begin{minipage}{0.92\linewidth}
\small
\textbf{Contract induction protocol}
\begin{enumerate}
  \item \textbf{Propose.} \code{propose\_contract} receives the instruction,
        active schemas, and policy. The configured builder model emits a candidate
        object, never a completion judgment.
  \item \textbf{Validate.} \code{validate\_contract} applies JSON-Schema validation,
        canonical goal IDs, exact-duplicate removal, policy checks, dependency-cycle
        detection, and adapter-supported verifier checks deterministically.
  \item If required fields are absent or the object is invalid, deterministic
        normalization runs first, followed by at most one schema-constrained repair
        call. Remaining uncertainty is written to $U$ rather than guessed.
  \item \textbf{Reduce.} \code{reduce\_contract} is the only operation allowed to
        change authoritative state in response to a validated contract event, tool
        observation, user confirmation, or authorized readback.
\end{enumerate}
\end{minipage}
\end{center}

In matched runs, the builder and any repair call use the same endpoint and
non-thinking setting as the main agent. Their calls, latency, input tokens, and
output tokens are charged to runtime overhead. AppWorld and $\tau^2$-bench use the
same output Schema and reducer semantics; adapters contribute only benchmark
vocabulary, policy clauses, entity types, and supported evidence recipes. Evidence
requirements are accepted only when a deterministic adapter maps the target tool
effect or policy clause to a readback, state-difference, receipt, or confirmation
predicate. Free-form model assertions are never admissible evidence.

Implicit omissions can be surfaced by contract validation, a tool precondition,
a user correction, or a Completion Gate deficit---never by consulting the benchmark
evaluator online. A material user update invokes \code{propose\_contract} only for
the affected subgraph; ordinary observations use \code{reduce\_contract} without
another model call. Previously verified effects remain in $L_t$ even when future
work is replanned. Every proposal, repair, rejection, and accepted contract edit is
versioned in the Trace Recorder.

Schema validity is not semantic adequacy. Development-set contract evaluation must
therefore annotate and report Goal Recall, Constraint Recall, Spurious Goal Rate,
Verifier Coverage, Contract Repair Rate, and Contract Failure Rate. These labels
and measurements are not yet available for the historical diagnostics reported
here; they are preregistered requirements for a confirmatory evaluation, not
observed evidence of contract quality.

\subsection{Three decision layers and completion}

We distinguish three outcomes throughout: (i) a model-proposed action or completion,
(ii) a runtime-authorized execution or gate-accepted completion, and (iii) an
evaluator-confirmed result over persisted state. A model's \code{complete\_task}
call is therefore an intent signal, not proof.

The Completion Gate can inspect only the observable projection $\widehat W_t$
constructed from authorized reads, not the latent world state $W_t$. It computes
\begin{equation}
\operatorname{AllowComplete}(H_t,E_t,L_t,P,\widehat W_t)=
\mathbb{I}\!\left[
\begin{aligned}
&\forall g\in G_{\mathrm{req}},\ \operatorname{Satisfied}(g,E_t,\widehat W_t)\\
&\land\ \forall \ell\in L_t^{\mathrm{intended}},\ \operatorname{Verified}(\ell,E_t,\widehat W_t)\\
&\land\ \neg\operatorname{UnresolvedUnexpectedEffect}(L_t^{\mathrm{unexpected}})\\
&\land\ \neg\operatorname{UnresolvedBlocker}(H_t)\\
&\land\ \operatorname{PolicySatisfied}(P,H_t,E_t)
\end{aligned}
\right].
\label{eq:completion}
\end{equation}

$G_{\mathrm{req}}$ contains validated goals marked required and not subsequently
withdrawn by the user. The ledger distinguishes $L_t^{\mathrm{intended}}$, effects
explicitly required by the contract; $L_t^{\mathrm{observed}}$, all effects actually
observed; and $L_t^{\mathrm{unexpected}}=L_t^{\mathrm{observed}}\setminus
L_t^{\mathrm{intended}}$. A non-idempotent or policy-sensitive write is always
recorded and verified, but it becomes a required effect only when it is contractually
intended. An unexpected effect blocks completion until it is rolled back or resolved
under an explicit policy. Thus, verifying that an unintended deletion occurred can
never convert it into a successful required effect. \textbf{Admissible evidence} comes from an authorized
tool, scoped user confirmation, or adapter-defined observation and is linked to one
contract item. \textbf{Freshness} uses an environment state version when available;
otherwise the adapter supplies a policy-specific TTL, and any relevant write
invalidates older evidence. Evidence conflicts are ordered by state version and
source authority, not recency alone. Incomparable active conflicts conservatively
block completion.

For an effect without a readback API, the adapter must define an alternative proof
such as an authenticated receipt plus a non-conflicting secondary observation or
explicit confirmation. Without such a recipe, the effect remains unresolved.
$\operatorname{Satisfied}$ and $\operatorname{Verified}$ are deterministic reducer or
adapter predicates; an LLM may extract candidate fields but cannot decide either
predicate. An \textbf{unresolved blocker} is a missing entity, confirmation,
permission, dependency, active conflict, or failed postcondition. A \textbf{false
acceptance} occurs when the gate allows completion but the offline evaluator fails.
A \textbf{false block} occurs when an offline snapshot evaluation shows that the
contract was already satisfied at a rejected attempt. The runtime never calls the
benchmark evaluator online.

\section{System Design}

\cref{fig:architecture} shows the complete control loop and the four top-level
modules. Solid arrows represent control flow; dashed lines represent trace events.
\begin{center}
\includegraphics[width=\linewidth]{figures/system_architecture.pdf}
\captionof{figure}{\safedesk{} runtime architecture. Context Manager constructs
each model call from runtime state and trace references. Completion Gate and
Response Grounding are subcomponents of State \& Verification, not additional
top-level modules. Trace Recorder receives events from every module.}
\label{fig:architecture}
\end{center}

\subsection{Runtime invariants and enforcement assumptions}
\label{sec:invariants}

The implementation is designed around three runtime invariants rather than a
claim to invent explicit state management. First, a verified non-idempotent effect
cannot be executed again with the same normalized effect fingerprint unless the
ledger records rollback evidence. Second, a completion request cannot be accepted
while an intended effect is unverified or an observed unexpected effect remains
unresolved. Third, the model cannot mutate authoritative runtime state: it can only
propose typed actions, and the State Reducer accepts only validated contract,
tool-observation, verification, or confirmation events. The first two statements
hold provided that all effectful tools traverse the Runtime Hook and adapters map
all relevant writes to effects; the third holds provided that the state-store
interface is not exposed to the model process. These are implementation invariants,
not claims that an adapter can observe an unobservable external side effect. The
prototype tests cover selected enforcement paths; they are not a formal proof that
an adapter has mapped every external write.

\subsection{State \& Verification}

State \& Verification implements the chain
\emph{Task Contract $\rightarrow$ Task State $\rightarrow$ Evidence Board
$\rightarrow$ Effect Ledger $\rightarrow$ Post-action Verification
$\rightarrow$ Completion Gate $\rightarrow$ Response Grounding}.

\paragraph{Task State.}
$H_t$ maintains subgoals, dependencies, constraints, resolved entities, blockers,
and execution phase. Reducer transitions are deterministic over typed events; the
model cannot directly mutate $H_t$.

\paragraph{Evidence Board.}
Each evidence item records source, timestamp, scope, validity, linked subgoal, and
whether later evidence supersedes it. Conflicts are retained for audit rather than
silently overwritten.

\paragraph{Effect Ledger.}
Each write records its intent, target entity, normalized parameters, idempotency
key, execution status, and verification status. This separates an attempted write,
a tool-reported success, and an observed environment effect.

\paragraph{Post-action Verification and completion.}
Verification prefers environment readback, object queries, state differences, and
tool-result consistency checks. Completion Gate never returns only ``incomplete'';
it emits missing subgoals, stale or conflicting evidence, unverified effects,
blockers, and unsatisfied policy conditions. Response Grounding applies the same
evidence boundary to the final natural-language response.

\subsection{Tool Execution Guard}

The guard follows:
\begin{center}
\small Model Tool Call $\rightarrow$ Action IR $\rightarrow$ Schema Validation
$\rightarrow$ Deterministic Repair $\rightarrow$ Policy/Precondition Check
$\rightarrow$ Dependency Scheduling $\rightarrow$ Effect Preflight
$\rightarrow$ Execute / Defer / Reject.
\end{center}

Action IR normalizes tool names, arguments, call identifiers, dependencies, and
effect class. Deterministic repair handles unambiguous aliases, scalar coercions,
and default-free formatting fixes; uncertain changes return a structured error to
the model. Policy checks enforce authentication, confirmation, permission, and
domain rules. Independent reads execute in parallel; dependent or state-changing
calls execute serially. Non-idempotent writes are serial by default.

\textbf{Defer never means discard.} We use explicit rejection and reproposal:
the current call receives a structured \code{deferred/not\_executed} result and
is permanently closed; after observing refreshed state, the model must issue a
new call with a new \code{tool\_call\_id}. No identifier is executed after a tool
result has been returned for it.

Deterministic rebinding is allowed only when the model explicitly supplies an
Action-IR placeholder such as \code{\$result\_1.contact\_id} and declares the
dependency edge. The runtime substitutes the unique observed field and revalidates
the call. It never replaces a model-guessed entity ID, chooses among multiple
candidates, or resolves semantic ambiguity; those cases return to the model or
user. Effect preflight compares normalized intent and idempotency fingerprints
against $L_t$. A bounded resolver may add an out-of-schema tool only in the
separate resolver augmentation described in \cref{sec:appworld-protocol}; the
matched core experiment disables expansion.

\subsection{Recovery Controller}

Recovery follows \emph{Failure Classification $\rightarrow$ Recoverability
Assessment $\rightarrow$ Typed Recovery $\rightarrow$ Progress Check
$\rightarrow$ Stagnation Detection $\rightarrow$ Local Replanning
$\rightarrow$ Safe Stop}. Failure classes include schema, parameter, entity,
permission, policy, transient tool, environment mismatch, duplicate action, no
progress, context loss, and unknown.

The controller deterministically repairs parameters when the correction is unique,
re-queries missing entities, applies bounded retry to transient faults, and reads
the environment before repairing a write mismatch. No-progress failures trigger
local replanning over the affected subgoal rather than a full restart. Permission
and policy failures cannot be bypassed. A non-idempotent write is never replayed
without evidence that the first attempt had no effect. Recovery budgets constrain
attempts per failure class and total turns. Progress compares changes in goals,
evidence, effects, and blockers; repeated equivalent reads or failures trigger a
safe stop with an explicit unresolved state.

Recovery success is typed rather than inferred from a later task result. A schema
recovery succeeds when the repaired call validates and executes; entity recovery
succeeds when re-query produces a uniquely valid entity; mismatch recovery succeeds
when a post-repair readback satisfies the expected predicate; and no-progress
recovery succeeds only when it adds admissible evidence or completes a new subgoal.
Task success after recovery remains a separate evaluator-level outcome.

\subsection{Context Manager}

Context Manager constructs four layers before each model call: (i) the original
task contract and hard constraints; (ii) current structured state; (iii) recent
complete interaction; and (iv) summaries of older history with references to raw
results. It must preserve entity identifiers, amounts, times, user confirmations,
prohibitions, tool dependencies, idempotency keys, and unresolved errors.

The token budget is partitioned among system instruction, task state, active tool
schemas, recent turns, evidence, and recovery context. Large observations are
projected to required fields, but their trace references remain available for
targeted retrieval. Compression invariants detect dropped identities, numeric
constraints, confirmations, dependencies, and active failures before the context
is sent to the model.

\cref{tab:trace-example} instantiates the four modules on the motivating example.
It shows why a plausible tool response, a verified environmental effect, and an
accepted completion are distinct runtime events.
\begin{table}[t]
\centering
\caption{Illustrative end-to-end trace for the calendar-and-notification example.
The example explains runtime semantics and is not counted as benchmark evidence.}
\label{tab:trace-example}
\small
\begin{tabularx}{\linewidth}{@{}cYYYY@{}}
\toprule
Step & Model proposal & Runtime module(s) & Runtime disposition & State, evidence, or effect update \\
\midrule
1 & Create event with missing time & Tool Execution Guard & Reject; return typed deficit & Required field remains unresolved \\
2 & Create event with complete arguments & Tool Execution Guard; State \& Verification & Execute serially & Planned effect is recorded; readback verifies event ID \\
3 & Repeat the same creation & Tool Execution Guard; Effect Ledger & Block duplicate write & Verified prior effect and idempotency key are cited \\
4 & Propose completion & Completion Gate & Gate blocks & Notification goal lacks admissible evidence \\
5 & Send notification & Tool Execution Guard; Post-action Verification & Execute, then verify & Notification effect becomes verified \\
6 & Propose completion & Completion Gate; Response Grounding & Gate accepts & All required goals, effects, and policies are satisfied \\
\bottomrule
\end{tabularx}
\end{table}


\section{Implementation}

\safedesk{} exposes a Runtime Hook around the harness model/tool loop. Adapters map
benchmark calls into Action IR and map observations into typed reducer events. The
State Reducer is the only component permitted to mutate runtime state. A persistent
store commits task state, trace events, and evaluator inputs atomically at each
turn. Configuration switches enable modules individually; disabling \safedesk{}
passes proposals to the original baseline path without changing the prompt or
evaluator.

\begin{center}
\begin{minipage}{0.88\linewidth}
\small
\hrule\vspace{4pt}
\textbf{Runtime main loop}
\begin{enumerate}[leftmargin=1.8em,itemsep=0pt]
  \item initialize contract and runtime state;
  \item while budget remains: build bounded context and obtain a model proposal;
  \item normalize actions; guard and schedule calls; execute authorized tools;
  \item update evidence and effects through the Reducer; verify writes;
  \item classify failures, apply bounded recovery, and detect stagnation;
  \item evaluate completion; persist state before evaluator access.
\end{enumerate}
\vspace{2pt}\hrule
\end{minipage}
\end{center}

Every tool action passes through the Runtime Hook, and every disposition receives
a protocol-compliant result. Trace records include proposals, Action IR, execute/
defer/reject decisions, tool observations, state transitions, verification, and
recovery. Redaction removes API keys, credentials, personal fields, and configured
sensitive arguments before persistence.

\section{Proposed Evaluation Protocol}

\subsection{AppWorld matched comparison}
\label{sec:appworld-protocol}

No independently administered held-out AppWorld split exists in the present
artifact. This subsection is therefore a preregistered construction protocol, not
a completed confirmatory experiment. Once an independent curator creates and
withholds a task split, the comparison will be Plain Baseline versus SafeDesk-core.
Both must
use the provider API model identifier \AWModelID{} through the OpenAI-compatible
endpoint \AWAPIEndpoint{}, with \code{thinking.type=disabled} and temperature
zero. The official documentation identifies \AWModelID{} as a provider API model
identifier supporting function calling and explicit thinking control
\citep{deepseek2026v4api,deepseek2026thinking}; it is not evidence of an immutable
backend snapshot. The historical baseline was accessed on \AWModelAccessDate, and
its response payloads did not retain a system fingerprint or provider build ID.
It is therefore a diagnostic observation only and is not a matched comparator for
any future SafeDesk result.

The official API Predictor runs once per task, and its ordered output of at most
\AWAPILimit{} APIs is persisted. Plain Baseline, the budget-matched control, and
SafeDesk-core replay the exact same per-task schemas. Dynamic Tool Resolver is
disabled in the primary comparison. SafeDesk+Resolver is a separate augmentation
that reports APIs added, schema tokens added, resolver calls, and tasks rescued;
it is never pooled with the core effect. DeepSeek V4 Pro is a reference condition,
not evidence about SafeDesk unless every remaining setting is matched.

The conditions are listed in \cref{tab:appworld-conditions}.
\begin{center}
\begin{minipage}{0.96\linewidth}
\centering
\captionof{table}{Preregistered AppWorld conditions. The held-out task set is not yet constructed.}
\label{tab:appworld-conditions}
\small
\begin{tabularx}{\linewidth}{@{}lYYr@{}}
\toprule
Condition & Model & Runtime & Tasks \\
\midrule
Plain Baseline & \AWModelID & Function calling; fixed schemas & future held-out \\
Resource control & \AWModelID & Plain agent; prescribed self-check schedule & future held-out \\
SafeDesk-core & \AWModelID & Four modules; fixed schemas & future held-out \\
SafeDesk+Resolver & \AWModelID & Separate dynamic-tool augmentation & future held-out \\
Strong reference & \texttt{deepseek-v4-pro} & Plain function calling & future held-out \\
\bottomrule
\end{tabularx}
\end{minipage}
\end{center}

For a confirmatory rerun, each task is assigned a deterministic seed and an
interleaved condition order; Plain, resource control, and SafeDesk-core runs occur
within the same provider window. The release stores the requested and returned
model fields, system fingerprint when supplied, endpoint, request schema hash,
base-policy hash, task-instruction hash, runtime context-renderer hash, Predictor
output, \AWAPILimit-API cap, a \AWMaxTurns{}-turn cap, persistence logic, evaluator,
retry policy, parallel-call setting, and token-accounting scope. ``Same prompt''
means the same base policy and task instruction; runtime-generated state and
evidence context are deliberately distinct and are versioned separately.
Independent reads may be parallelized by the runtime, but both conditions receive
the same model-level parallel-tool setting.

\subsection{Budget accounting and controls}

Resource accounting distinguishes: (i) \emph{agent turns}, model proposal cycles
initiated by the harness; (ii) \emph{total LLM calls}, including Predictor,
Contract Builder, repair, classifier, replanning, and summarization calls;
(iii) \emph{agent-proposed tool calls}; (iv) \emph{runtime-generated verification
calls}; (v) \emph{total executed tool calls}; (vi) \emph{total tokens} over every
LLM call; and (vii) wall-clock latency. Verification calls consume the tool-call
budget. Builder, classifier, recovery, resolver, and summarizer tokens contribute
to total tokens. A recovery model call is not an agent turn but is a total LLM
call. Runtime-assisted calls use the same model ID and thinking setting as the
agent; no stronger helper is permitted.

The primary protocol fixes a per-task vector of total LLM input and
output tokens, total LLM calls, executed tool calls, and environment reads. Every
condition stops when any component is exhausted; verification reads count against
the same environment-read and tool-call budgets. The resource-control condition is
not matched post hoc. It receives exactly three predeclared controls: (A) one fixed
self-check after every write, (B) one fixed reflection after every five agent
turns, and (C) one fixed self-verification after every model-proposed
\code{complete\_task}. It has no structured Task State, Evidence Board, Effect
Ledger, Guard, or completion logic. All conditions share the same token, call,
tool, and read ceilings. A secondary resource--performance analysis repeats the
study at predeclared token and tool-call budgets. Once the \AWMaxTurns{} agent-turn cap is
reached, a condition may only persist state and emit a safe-stop record. No
efficiency or reliability conclusion is made until it survives these controls.

\subsection{Data validity and runner repairs}

The AppWorld runner originally evaluated some tasks before flushing the in-memory
world to the output database. The corrected runner persists through AppWorld's
state-save API before every evaluation and records a per-task persistence flag;
all \AWTaskCount{} final rows satisfy that flag. Separately,
\AWGBKRepairCount{} Windows GBK infrastructure failures were rerun with
interpreter-level UTF-8, and \AWGBKRepairSuccessCount{} produced evaluator success
after the infrastructure repair. Only these infrastructure failures were replaced;
model or evaluator failures were not selectively rerun. The same persistence and
encoding fixes apply to every future condition.

All \AWShardCount{} shards were re-evaluated offline after persistence. The canonical sorted
task-list and repair-subset SHA-256 values are reported in
\cref{app:reproducibility}. Final infrastructure errors equal
\AWFinalInfraErrorCount{}. \AWInvalidDurationCount{} historical negative duration values caused by
environment time freezing are excluded; future runs use a monotonic host clock.
Thus ``validated baseline'' below means a frozen task list with runner-level
infrastructure repairs, not a post-hoc model correction.

\subsection{Test-set exposure and confirmatory isolation}

The historical AppWorld \code{test\_challenge} traces and task-level diagnostics
were inspected while developing this runtime. The current \AWTaskCount{}-task
aggregate is therefore explicitly exploratory and is not used for confirmatory
statistics, model selection, or a claimed SafeDesk gain. A pre-exposure design
freeze cannot be reconstructed from the available workspace history. Confirmatory
evaluation requires a newly constructed split held by an independent curator or
evaluation service. Before access, the release must archive the source
revision, configuration, base-policy hash, context-renderer hash, thresholds,
budget vector, task-order generator, and analysis script. This is an experimental
boundary, not a post-hoc disclaimer.

\subsection{\texorpdfstring{$\tau^2$-bench}{tau2-bench} matched comparison}

The future comparison pins the local $\tau^2$-bench checkout to commit
\TauBenchmarkRevision{} (package \TauPackageVersion{}), uses the recommended
\code{base} split and text-based half-duplex setting, and records the sorted
task-list SHA-256 \TauTaskListHash{}. The current checkout's task-fix status is
``not independently audited''; later maintenance fixes are excluded unless the
pin and hash are deliberately updated. The agent, user simulator, and judge model
identifiers remain unassigned in this manuscript, so no $\tau^2$ result is claimed.
Before the run, all three must be frozen in a manifest. The formal comparison then
uses the selected agent with and without \safedesk{}, with GLM-5.2 as a reference
only under an otherwise matched protocol.
Z.AI documents \texttt{glm-5} function calling and explicit thinking control, while
the GLM-5.2 model card is cited only for the reference identity
\citep{zai2026glm5,zai2026glm52,barres2025tau2,tau2benchsoftware2026}. The historical
endpoint was \TauAgentEndpoint{} on \TauModelAccessDate; its provider revision was
\TauProviderVersion and cannot be reconstructed. A formal run records the
evaluator command, version and configuration, agent/user/judge response metadata,
seeds, turn limits, tool limits, and all role-separated token counts. Airline,
retail, and telecom
tasks each run \TauTrialsPerTask{} times, and the official saved-result procedure
computes Pass$^1$ and Pass$^4$. We report per-domain and macro scores, policy
violations, premature tool calls, response inconsistencies, turns, and agent-only
tokens. The available single-run diagnostics are retained only as local adapter
smoke evidence and are not reported as benchmark performance.

\subsection{Metrics and statistical analysis}
\label{sec:metrics}

Metric definitions and denominators appear in \cref{app:metrics}. The primary TGC
inference is a paired scenario-clustered bootstrap, with a clustered permutation
test as a robustness check. On the future split, SGC treats its independently
reported scenario count directly; the historical \AWScenarioCount{} count is used
only in the appendix.
The task transition matrix and ordinary McNemar test are auxiliary task-level
descriptions because three variants within a scenario are correlated. Difficulty
intervals also resample scenarios rather than individual tasks. Token, call, and
turn differences use \AWBootstrapSamples{} paired bootstrap samples clustered by
scenario; Holm correction controls multiplicity across ablations. P-values,
paired contrasts, and transition counts are withheld until emitted by the frozen
statistics script; no value is inferred from aggregate macros.

\section{Prespecified Ablation Design}

The primary necessity analysis removes State \& Verification, Tool Execution
Guard, Recovery Controller, or Context Manager one at a time from the full runtime.
The linked outcomes are respectively completion/false block, invalid and duplicate
writes, stagnation/rescue, and token-tail statistics with TGC. A secondary
incremental sequence adds the modules in that order for engineering interpretation.
Incremental gains depend on module order and cannot be interpreted as independent
causal contributions. The versioned configuration and analysis script must be
released before held-out evaluation; no ablation outcome is claimed here.

\section{Discussion}

\paragraph{Why a runtime is useful.}
The model proposes candidate actions; the runtime maintains state consistency,
permissions, idempotency, verification, audit, and budgets. This division makes
non-negotiable controls explicit and reusable across model versions.
It also creates a stable experimental boundary: proposal quality can change while
the execution contract remains fixed. This does not make the runtime infallible;
it makes its decisions inspectable, replayable, and separately testable from the
language model that produced a candidate action.

\paragraph{Relationship to agentic RL.}
RL can improve the proposal distribution and long-horizon policy, while \safedesk{}
governs the consequences of sampled proposals. A trained agent can therefore run
inside the same runtime; the approaches are complementary.
The matched experiment deliberately holds model weights fixed so that the runtime
effect is identifiable. A subsequent factorial study can cross training and
runtime conditions to test whether better proposals reduce intervention frequency
and whether governance preserves gains under rare but consequential failures.

\paragraph{Completion Gate conservatism.}
An incomplete Evidence Board can block an already completed task. Readback adds
latency and tokens, and an overly strict gate can reduce TGC. False-block rate,
verification coverage, and governance overhead must therefore accompany false
completion and TGC.
The appropriate threshold is domain dependent: a redundant read may be justified
for a payment or account change but excessive for a reversible search preference.
The gate therefore consumes evidence requirements from the Task Contract instead
of applying one universal proof rule to every action.

\paragraph{Dynamic tool expansion.}
Resolver expansion can recover a missing API, but also increases schema tokens,
ambiguity, attack surface, and the probability of selecting the wrong tool.
Expansion must be scoped to the active subgoal, bounded, and policy checked.
Candidate APIs should be logged with their retrieval reasons and rejected
alternatives, so a capability gain is not confused with an unreported expansion
of the action space. Matched evaluation must expose the same underlying catalog
even when active subsets change over time.

\paragraph{Context compression risk.}
Summaries can drop entity identifiers, user confirmations, numeric constraints,
tool dependencies, or prohibitions. Context invariants and trace-backed retrieval
reduce this risk but cannot eliminate imperfect summarization.
SafeDesk therefore treats the structured contract and current state as canonical,
not as another free-form summary. Raw observations remain addressable by trace
reference, allowing targeted recovery without repeatedly placing the full history
in the model context.

\paragraph{Module interaction.}
Recovery depends on accurate task and failure state; completion depends on
admissible evidence; context construction determines what recovery information the
model can see; and Tool Execution Guard affects both side effects and stagnation.
Accordingly, leave-one-out and interaction analyses are more informative than an
incremental sequence alone.
These dependencies also explain why module-local metrics cannot substitute for
task success. Schema repair can reduce invalid calls while selecting the wrong
entity, and aggressive compression can reduce tokens while suppressing evidence
needed by the gate. Mechanism metrics must therefore be paired with
evaluator-confirmed outcomes.

\section{Limitations}

\begin{enumerate}
  \item Contract induction can split a goal incorrectly or omit an implicit subgoal.
  \item Completion Gate is only as reliable as captured and fresh evidence.
  \item Some external writes lack a readback API or observable state difference.
  \item Idempotency fingerprints can conflate semantically distinct actions with similar arguments.
  \item Dynamic Tool Resolver can expand the tool boundary and attack surface.
  \item Context compression can discard a critical fact despite invariant checks.
  \item Provider aliases and model behavior can drift unless immutable versions are recorded.
  \item AppWorld and $\tau^2$-bench do not represent all production domains or policies.
  \item Runtime checks add engineering complexity, latency, storage, and failure modes.
  \item Domain-specific policies, evidence extractors, and verification adapters still require engineering.
\end{enumerate}

\section{Ethics and Safety}

\safedesk{} is designed to reduce uncontrolled side effects, not expand agent
privilege. Deployment should apply least privilege, explicit confirmation,
sensitive-field redaction, bounded retries, rollback where available, and human
takeover. Dynamic resolution must not bypass policy. Benchmark tools run in
isolated environments; released traces must exclude keys, credentials, personal
data, and replayable account state.

\section{Artifact Availability}

The local artifact contains the runtime core, versioned JSON Schemas, prototype
tests, adapters, trace format, figure scripts, and paper sources. It is \emph{not}
currently a public archival release: there is no public repository URL, release
commit, archival test log, or assigned artifact license. The workspace contains a
dependency lock and local tests, but readers cannot independently reproduce the
reported prototype results until an archival release is made. A public version must
pin source revision and dependencies, publish test logs and mock definitions,
redact secrets from traces, state its license and release date, and include the
frozen evaluation configuration required by \cref{app:reproducibility}.

\section{Conclusion}

\safedesk{} reframes long-horizon tool execution from a free-form message loop as
state- and evidence-driven runtime governance. Its four modules govern completion
judgment, tool execution, failure recovery, and context efficiency, while preserving
the distinction between model proposals, runtime authorization, and evaluator
outcomes. Prototype validation demonstrates the control path, while matched
SafeDesk, repeated-trial $\tau^2$-bench, and ablation effects remain unmeasured
and must not be inferred from historical diagnostics or synthetic targets. Runtime
governance complements stronger models and agentic training rather than replacing them.

\appendix
\section{Metric Definitions}
\label{app:metrics}

This appendix fixes the unit of analysis and denominator for every reported
quantity. A task is valid only when the benchmark environment, persistence step,
and evaluator complete without an infrastructure error. Scenario-level measures
use only complete AppWorld scenario triplets; call-level measures use either all
model proposals or executed actions as stated below. Rejected proposals therefore
remain visible to validity metrics but cannot contribute an environmental side
effect. Rates are computed from counts before rounding and are displayed to two
decimal places. Absolute differences between rates are reported in percentage
points, whereas relative changes are reserved for quantities such as token use.

Completion metrics preserve temporal and causal scope. A proposal records model
intent, a gate decision records runtime authorization, and evaluator success is a
judgment over the persisted final state. One task may make several attempts, be
blocked, gather additional evidence, and later be accepted; attempt sets and task
sets are consequently not interchangeable. Recovery follows the same rule: repair
of a classified failure is a recovery-attempt outcome, while eventual benchmark
success remains a task outcome. These definitions are shared by the tables,
figures, trace exporter, and statistics scripts.

\begin{center}
\begin{minipage}{\linewidth}
\centering
\captionof{table}{Completion and recovery metrics, units, and denominators.}
\label{tab:completion-metric-definitions}
\small
\begin{tabularx}{\linewidth}{@{}lYY@{}}
\toprule
Metric & Definition & Unit / denominator \\
\midrule
TGC & Tasks passing every evaluator test divided by valid tasks & task \\
SGC & Complete scenarios for which all three variants pass divided by complete scenarios & scenario \\
System-accepted completion failure & Evaluator-failed tasks after the system allowed termination divided by tasks whose termination the system allowed & accepted task \\
Model-proposed completion failure & Evaluator-failed tasks with a model completion proposal divided by tasks with such a proposal & proposal task \\
False-block attempt rate & Rejected attempts whose decision-time snapshot already satisfied the contract divided by all rejected attempts & blocked attempt \\
False-block affected tasks & Tasks containing at least one false-block attempt divided by tasks blocked at least once & blocked task \\
Gate rescue rate & Tasks blocked at least once and later successful divided by tasks blocked at least once & blocked task \\
Blocks per rescued task & Blocked attempts among rescued tasks divided by rescued tasks & rescued task \\
Max-turn rate & Tasks reaching the configured turn cap divided by valid tasks & task \\
Verification coverage & Required effects with admissible verification divided by required effects & required effect \\
Recovery-attempt success & Recovery attempts that repair their classified failure divided by triggered recovery attempts & recovery attempt \\
Task success after recovery & Recovered tasks finally passing the evaluator divided by tasks that triggered recovery & task \\
\bottomrule
\end{tabularx}
\end{minipage}
\end{center}

\begin{center}
\begin{minipage}{\linewidth}
\centering
\captionof{table}{Tool reliability and resource metrics, units, and denominators.}
\label{tab:tool-resource-metric-definitions}
\small
\begin{tabularx}{\linewidth}{@{}lYY@{}}
\toprule
Metric & Definition & Unit / denominator \\
\midrule
Invalid-call rate & Invalid proposed calls divided by all model-proposed calls & proposed call \\
Out-of-schema rate & Proposed calls absent from the active schema divided by all model-proposed calls & proposed call \\
Duplicate-call rate & Semantically duplicate executed calls divided by all executed calls & executed call \\
Duplicate-write rate & Duplicate write actions divided by all executed writes & executed write \\
Token statistics & Mean, median, P95, and maximum, separated by agent, predictor, user simulator, judge, and runtime & token / task \\
Latency statistics & Mean, median, P95, and maximum from induction to persistence or safe stop & second / task \\
\bottomrule
\end{tabularx}
\end{minipage}
\end{center}

\paragraph{Completion scopes.}
A \emph{model-proposed completion} is an action proposal and does not imply that
the runtime accepted termination. A \emph{gate-accepted completion} is a runtime
decision based on the evidence available at that step. An
\emph{evaluator-confirmed success} is the benchmark's judgment over the persisted
final state. The three events are recorded independently. Multiple completion
attempts can belong to one task, so attempt-level and task-level denominators are
never interchanged.

For the plain baseline, every \code{complete\_task} proposal is implicitly accepted,
so model-proposed and system-accepted completion-failure rates coincide. For
SafeDesk, blocked proposals remain in the model-level rate but not the
system-accepted denominator. A false-block label is produced only after the run:
the evaluator receives the immutable, redacted database snapshot and serialized
Task Contract captured immediately before the rejected attempt, executes the
standard offline test suite without the agent or runtime, and records whether all
contract-required benchmark predicates were already satisfied. Snapshots that fail
to restore, lack evaluator coverage, or include an unresolved policy confirmation
are excluded and reported separately; they are never silently counted as either
true or false blocks. No evaluator output is available to the online gate.

\paragraph{Call relations.}
Out-of-schema calls are a subset of invalid calls and are not added to invalid
calls. Invalid-call rates use all model-proposed calls, including rejected calls.
Duplicate-call rates use executed calls because a rejected proposal cannot repeat
an environmental effect. An executed call is duplicate only when it has the same
normalized tool name and equivalent normalized arguments as an earlier call, no
intervening relevant state or evidence version changed, and the call is neither
pagination nor an explicit readback verification. Duplicate-write rates use
executed writes. A write is duplicate only when intended effect, target entity,
and critical normalized arguments match and the first write was already observed
or verified. These predicates are computed from Action IR and reducer versions,
not assigned by free-form trace review.

\paragraph{Typed recovery outcomes.}
Recovery success is evaluated against the failure-class exit predicate: schema
recovery requires a validating, executable call; entity recovery requires a unique
valid entity after re-query; environment-mismatch recovery requires readback to
satisfy the expected postcondition; and no-progress recovery requires new
admissible evidence or completion of a previously incomplete subgoal. Repairing
this local predicate does not imply evaluator-confirmed task success.

\paragraph{Resource accounting.}
Agent turns count harness proposal cycles. Total LLM calls additionally include
Predictor, Contract Builder, repair, classifier, replanning, resolver, and
summarizer calls. Agent-proposed tool calls and runtime-generated verification
calls are separate proposal sources; total executed calls counts both after guard
disposition. Total tokens aggregate every model role and are also reported by
role. Wall-clock latency spans contract induction through final persistence or
safe stop.

\paragraph{$\tau^2$-bench.}
Pass$^1$ is the fraction of tasks succeeding in a specified trial. Pass$^4$ is
the fraction succeeding in all \TauTrialsPerTask{} fixed repeated trials. Policy
violation is a tool action that violates a domain rule. Premature tool call is an
action taken before required identity, confirmation, or policy preconditions are
met. Tool-response inconsistency means the final response contradicts observed
tool state. Average turns and agent-only tokens exclude user-simulator and judge
usage.

\section{Reproducibility Checklist}
\label{app:reproducibility}

\begin{itemize}[leftmargin=*,itemsep=2pt]
  \item Freeze benchmark commit, task list, task order, evaluator, and database seed.
  \item Record exact provider-returned model identifier, non-thinking mode,
        temperature, prompt hash, schema hash, and API Predictor output.
  \item Use unique experiment, task, trial, and configuration identifiers.
  \item Persist the final environment state before invoking the evaluator.
  \item Save every model proposal, runtime disposition (execute, defer, reject),
        tool result, state reduction, verification, recovery, and completion decision.
  \item Close a deferred \texttt{tool\_call\_id} with one
        \texttt{deferred/not\_executed} result; any later proposal uses a new ID.
  \item Substitute dependent values only for explicit Action-IR placeholders with
        a declared edge and one unambiguous observed output; never repair guessed IDs.
  \item Separate agent, predictor, user-simulator, judge, and runtime token usage.
  \item Classify infrastructure failures and rerun only under the identical frozen condition.
  \item Redact API keys, credentials, personal information, and sensitive tool fields from traces.
  \item Run the result-macro validator before compiling the paper; projected-target
        macros are forbidden in manuscript sources.
  \item Report the code commit, environment lock hash, source-data hashes, and
        generated-result hashes with the final artifact.
\end{itemize}

\subsection{Frozen diagnostic configuration}

\begin{center}
\small
\begin{tabularx}{\linewidth}{@{}lYY@{}}
\toprule
Field & AppWorld diagnostic & $\tau^2$-bench diagnostic \\
\midrule
Agent model ID & \AWModelID & \TauAgentModelID \\
Endpoint & \AWAPIEndpoint & \TauAgentEndpoint \\
Access date & \AWModelAccessDate & \TauModelAccessDate \\
Provider revision & \AWProviderVersion & \TauProviderVersion \\
Benchmark package/revision & AppWorld revision not archived & $\tau^2$ \TauPackageVersion; \texttt{1901a301961c} (full hash in manifest) \\
Thinking & disabled & disabled \\
Temperature & \ModelTemperature & \ModelTemperature \\
\bottomrule
\end{tabularx}
\end{center}

The AppWorld runner used an API Predictor capped at \AWAPILimit{} active APIs and
an \AWMaxTurns{}-turn cap. Provider-returned immutable build identifiers were not
preserved in the historical runs, so the endpoint, model alias, and access dates
bound the result more precisely than an unsupported version claim. The historical
run is exploratory because its task-level traces informed development. Formal reruns
must persist the response model field, system fingerprint when available, and
deployment revision when supplied.

\subsection{Runner repair audit}

The historical runner repair log records a state-save correction and
\AWGBKRepairCount{} Windows GBK decoding failures rerun under interpreter-level
UTF-8, of which \AWGBKRepairSuccessCount{} subsequently passed. The original
runtime package version, environment lock, and repair execution logs were not
archived together, so these counts are not used as a performance result or a
reproducible repair-rate claim. Every future condition must record the persistence
event, environment lock, runner revision, and exact rerun reason; model failures
must never be selectively rerun. The available historical identifiers are:

\begingroup
\raggedright
\begin{itemize}[leftmargin=*,itemsep=2pt]
  \item canonical sorted AppWorld task list SHA-256:
        \texttt{\AWTaskListHashDisplay};
  \item UTF-8 repair subset SHA-256:
        \texttt{\AWRepairListHashDisplay}.
\end{itemize}
\endgroup

\AWInvalidDurationCount{} historical negative duration values produced by environment time freezing
are excluded from the archived duration summaries only. Future conditions use a
monotonic host clock and retain the raw timing stream.

\section{Implementation Inventory}
\label{app:implementation}

For artifact inspection only, the current workspace contains a framework-neutral
runtime interface, versioned contracts, a DeerFlow hook, one tested AppWorld
integration, and prototype $\tau^2$-bench adapters. Code size and schema counts
are intentionally excluded from the main argument: they describe artifact scope,
not correctness. The current workspace has
no recoverable pre-exposure Git history, so it cannot prove a historical design
freeze. The public release must publish a source revision, dependency lock, and
CI log before confirmatory benchmark access.

\begin{center}
\small
\begin{tabularx}{\linewidth}{@{}lY@{}}
\toprule
Scope & Local validation outcome and boundary \\
\midrule
Runtime and AppWorld integration & \texttt{python -m pytest} over configured paths: \LocalUnitTestPassed{} passed. Contract, state-store, guard, trace, and public-catalog tests passed locally. \\
End-to-end runtime prototype & \PrototypeValidationPassed{}/\PrototypeValidationRepetitionCount{} repetitions of one deterministic mock regression scenario passed; \PrototypeValidationCrashes{} crashes. Each used SQLite state and trace stores, rejected one invalid call, verified one write, blocked unverified completion, prevented a duplicate write, then reloaded durable state and trace. This is not an AppWorld score. \\
DeerFlow live compatibility & \LocalUnitTestSkipped{} skipped because \texttt{deerflow.sandbox.tools} is not installed. This integration is not verified. \\
AppWorld public catalog & The real API-document fixture yields \AppWorldCatalogToolCount{} tools. Schema conversion is unit-tested; this is not an end-to-end benchmark run. \\
$\tau^2$ adapter & Local smoke invocation only. No repeated-trial or official Pass$^k$ claim is made. \\
\bottomrule
\end{tabularx}
\end{center}

The test outcome above is a local validation record, not a substitute for an
external CI run, baseline-bypass equivalence test, or evaluator-confirmed agent
success. Those release checks remain required before a confirmatory-results paper.

\section{Executable Runtime Specification}
\label{app:runtime-specification}

This appendix specifies the control path used by the prototype. It is not a claim
that the following algorithms have been benchmarked on AppWorld or $\tau^2$-bench.

\paragraph{Main loop.} For task $i$, initialize a validated contract and durable
aggregate. While no budget is exhausted, build a bounded Context Pack; obtain a
model proposal; normalize each proposal into Action IR; validate schema, policy,
dependencies, and effect preflight; execute the allowed schedule; reduce tool
observations; verify state-changing effects; classify any failure; apply at most the
configured typed recovery; and assess progress and completion. Persist state and
trace records before another model call or offline evaluation. A rejected or
deferred call receives a structured result and is never silently discarded.

\paragraph{Contract Builder.} Given instruction $u$, active schemas $\Sigma$, and
policy $P$, produce candidate $Q$. Validate $Q$ against
\code{agentgate-core/src/agentgate\_core/schemas/v1/task-contract.schema.json};
canonicalize identifiers; reject dependency cycles and unsupported verifier
recipes; and place missing fields in \code{unresolved}. A deterministic normalizer
precedes at most one schema-constrained repair. The reducer records the accepted
contract version. The model cannot mark a goal completed in this procedure.

\paragraph{Completion Gate.} For every required subgoal, require
\code{completed\_verified}. For every required evidence and effect identifier,
require a fresh verified record. Reject on unresolved failure, pending confirmation
or approval, active evidence conflict, an \emph{active unresolved} unintended side
effect, or duplicate irreversible effect. An unexpected effect with status
\code{observed}, \code{rollback\_pending}, or \code{unresolved} blocks completion;
\code{rolled\_back} and \code{policy\_accepted} do not. Return every blocker with
its type, reference, and recommended
phase; accept only if the blocker set is empty.

\paragraph{Effect fingerprint.} For a state-changing Action IR $a$, canonicalize
the tool name, operation, resource type and identifier, and only the schema-marked
critical arguments. Serialize the canonical object with sorted keys and compute
\code{sha256}. The ledger uses this value as the idempotency key. A matching effect
that is verified yields \code{already\_applied}; an in-flight or unknown effect
yields \code{verify\_first}; a failed effect yields \code{recovery\_required}; and
a rolled-back effect may be reserved again. Effect reservation and aggregate
persistence use an optimistic compare-and-set revision. A competing reservation
that loses the idempotency or persistence race is not dispatched; it receives a
state-refresh/replan result. The implementation state machine is
\code{ABSENT $\rightarrow$ RESERVED $\rightarrow$ IN\_FLIGHT $\rightarrow$
APPLIED\_UNVERIFIED $\rightarrow$ VERIFIED}. A restored \code{IN\_FLIGHT} effect
becomes \code{UNKNOWN} and must follow \code{VERIFY\_FIRST}; \code{FAILED} follows
typed recovery; \code{VERIFIED $\rightarrow$ ROLLED\_BACK $\rightarrow$ RESERVED}
is the only replay path. This protocol is atomic at the persisted task-aggregate
revision, not a claim of an independently normalized SQL effect table.

\begin{center}
\small
\begin{tabularx}{\linewidth}{@{}p{0.23\linewidth}p{0.23\linewidth}YY@{}}
\toprule
Invariant & Enforcement location & Assumption & Local counterexample test \\
\midrule
Verified non-idempotent effect is not re-dispatched & Effect Ledger reservation and CAS persistence & every effectful tool is mapped to an effect & duplicate key; two SQLite store instances race \\
Completion has verified intended effects and no active unexpected effect & Completion Gate and verifier results & contract and verifier recipes are complete & missing effect; observed versus rolled-back/policy-accepted unexpected effect \\
Model cannot mutate authoritative state & typed State Reducer and store boundary & model process has no store capability & invalid direct transition and reducer transition tests \\
\bottomrule
\end{tabularx}
\captionof{table}{Runtime invariants, enforcement assumptions, and selected local
counterexample tests. This is executable test evidence, not a formal proof.}
\end{center}

\begin{center}
\small
\begin{tabularx}{0.94\linewidth}{@{}p{0.20\linewidth}YY@{}}
\toprule
Failure class & Typed recovery & Boundary \\
\midrule
Schema or parameter & deterministic normalization, then one constrained repair & no guessed fields \\
Entity & re-query and bind a unique identifier & ambiguous entities block \\
Transient tool & bounded retry after backoff & retry budget is finite \\
Environment mismatch & read back and reconcile expected effect & no blind write replay \\
Duplicate action & verify existing effect & non-idempotent writes do not replay \\
No progress or context loss & local replanning with state references & safe stop on repeated stagnation \\
Permission or policy & report blocker / request confirmation & runtime never bypasses policy \\
\bottomrule
\end{tabularx}
\captionof{table}{Failure-to-recovery mapping. Recovery success means the local
exit predicate is repaired; it does not imply evaluator-confirmed task success.}
\end{center}

\begin{center}
\small
\begin{tabular}{@{}lrp{0.53\linewidth}@{}}
\toprule
Context component & Budget priority & Must retain \\
\midrule
Task contract and hard constraints & fixed & goal IDs, prohibitions, confirmations \\
Current state and blockers & fixed & entity IDs, dependencies, unresolved failures \\
Tool schemas & bounded & active tools and required completion tool \\
Recent interactions & bounded & latest tool result and recovery context \\
Evidence and effects & bounded & verified effects, idempotency keys, freshness \\
Older history & references only & trace IDs and deterministic summaries \\
\bottomrule
\end{tabular}
\captionof{table}{Context Budget policy. The future experiment must freeze numeric
allocations in its run manifest; this paper does not claim a benchmark-optimized
allocation.}
\end{center}

\paragraph{Contract example.} The following abbreviated object illustrates the
versioned Task Contract schema. It exposes goals/subgoals, dependencies, hard and
policy constraints, allowed and forbidden effects, confirmation requirements, and
evidence requirements. Unresolved fields are deliberately represented in Task
State rather than the immutable contract, as specified in \cref{sec:contract-induction}.
\begin{verbatim}
{"task_id":"example-1", "version":1,
 "original_instruction":"Create and verify one quarterly record.",
 "normalized_goal":"Create one verified quarterly record.",
 "subgoals":[{"subgoal_id":"create", "description":"Create record",
 "dependency_ids":[], "completion_condition_ids":["record-verified"],
 "required":true}],
 "constraints":[{"constraint_id":"no-duplicates",
 "description":"Do not create a duplicate record", "kind":"hard"},
 {"constraint_id":"policy", "description":"Use catalog tools only",
 "kind":"policy"}],
 "completion_conditions":[{"condition_id":"record-verified",
 "description":"Readback verifies the record", "required_effect_ids":["effect-42"]}],
 "allowed_effects":["create_record"], "forbidden_effects":["delete_record"],
 "required_confirmations":[]}
\end{verbatim}

\section{Prototype-Test Coverage}
\label{app:prototype-tests}

The \PrototypeValidationRepetitionCount{} end-to-end mock executions are
repetitions of one deterministic regression scenario, not \PrototypeValidationRepetitionCount{}
independent agent tasks. Each repetition checks the same path: one invalid call,
one typed recovery, one write/readback, one blocked completion, one duplicate-write
suppression, and durable SQLite state/trace reload. We report this distinction to
avoid interpreting the count as benchmark diversity.

The local suite additionally contains the coverage categories in
\cref{tab:prototype-coverage}. These are component and fault-injection tests, not
evaluator-confirmed AppWorld tasks. The current suite has zero property-based tests;
the concurrency and restart cases are deterministic fault-injection tests. The
public release must include the full test log and environment lock before these
local results can be independently reproduced.

\begin{center}
\small
\begin{tabularx}{\linewidth}{@{}p{0.19\linewidth}p{0.30\linewidth}Y@{}}
\toprule
Coverage category & Cases exercised & Enforcement boundary \\
\midrule
Schema errors & unknown tool and invalid arguments & schema guard returns a structured non-execution result \\
Dependency scheduling & dependency disposition; read/write scheduling & dependent writes defer until prerequisites are observed \\
Effect Ledger & duplicate key, in-flight/unknown, rollback transition & reservation and status-transition rules \\
Completion Gate & missing goal/effect; five unexpected-effect resolution states & active unresolved effects block; rolled-back or policy-accepted effects do not \\
Recovery & classified failure, measurable progress, bounded result & local recovery success is separate from task success \\
Context & bounded Context Pack construction and state references & hard task state remains canonical \\
Persistence & SQLite close/reopen; in-flight becomes unknown on restore & restart requires verify-first rather than blind replay \\
Concurrency & two SQLite store instances race to persist an effect & optimistic revision conflict prevents a second dispatch \\
\bottomrule
\end{tabularx}
\captionof{table}{Local prototype-test coverage. Categories overlap and do not
constitute independent benchmark scenarios.}
\label{tab:prototype-coverage}
\end{center}

\begin{center}
\small
\begin{tabular}{@{}lr@{}}
\toprule
Quantity & Current local record \\
\midrule
End-to-end regression-scenario repetitions & \PrototypeValidationRepetitionCount{} \\
Coverage categories & 8 \\
Property-based tests & 0 \\
Deterministic restart/concurrency fault-injection categories & 2 \\
AppWorld evaluator-confirmed SafeDesk tasks & 0 \\
\bottomrule
\end{tabular}
\captionof{table}{Scope of the prototype evidence.}
\end{center}

\section{Historical Exploratory Diagnostics}
\label{app:historical-diagnostics}

This appendix isolates a historical AppWorld \code{test\_challenge} diagnostic
from the method and protocol claims. The plain function-calling condition completed
\AWBaseSuccessCount/\AWTaskCount{} tasks (TGC \AWBaseTGC) and
\AWBaseScenarioSuccessCount/\AWScenarioCount{} complete three-variant scenarios
(SGC \AWBaseSGC). SGC here is scenario-level variant robustness, not multi-run
stability. The traces and task-level diagnostics were inspected during development;
therefore these values are neither a held-out comparison nor evidence of a SafeDesk
effect. No SafeDesk condition was run against this aggregate.

The local $\tau^2$-bench adapter runs are likewise smoke diagnostics only. They do
not use a frozen agent/user/judge manifest, role-separated accounting, or four
trials per task, and are not reported as benchmark performance. The future protocol
in \cref{sec:appworld-protocol} and the $\tau^2$ subsection define the required
replacement evidence.


\ifdefempty{\AuthorContributions}{}{\section*{Author Contributions}
\AuthorContributions}
\ifdefempty{\ArtifactLicense}{}{\section*{Artifact License}
\ArtifactLicense}

\bibliographystyle{plainnat}
\bibliography{references}

\end{document}
